% GENERATED FILE. Edit canonical lessons, then run npm run generate:books.
% canonical-source-hash: fnv1a64:a8c675cb329e4eb1
% canonical-lessons: TE-C549-amarcu, TE-C549-paracu, TE-C549-veligincu, TE-C549-ciluku, TE-C549-alasipovu, TE-R549-review-542-549, TE-R549-review-final-549

\chapter{Making a Home}
\label{ch:te-making-a-home}
\hlchaptermodality{549}

\begin{chapteropening}
\textbf{By the end of this chapter:} \emph{I can say that we set up, spread, lit, churned and wore ourselves out.}

Complete the last lesson of chapter 549: I can say that we set up, spread, lit, churned and wore ourselves out.
\end{chapteropening}

\section[\texorpdfstring{\te{అమర్చు} (amarcu)}{అమర్చు (amarcu)}]{\te{అమర్చు} (amarcu) — to set up; to arrange}

\label{lesson:TE-C549-amarcu}



\begin{quote}
Before the new one: say the Telugu for to sift, then the Telugu for tempering.
\end{quote}

\subsection*{You'll want to know: \te{అమర్చు}}
\textbf{\te{అమర్చు}} — \emph{amarcu} — \textquotedblleft{}to set up; to arrange\textquotedblright{}.

\begin{quote}
\te{కొత్త} \te{ఇంట్లో} \te{మేము} \te{సామాను} \te{అమర్చాము}.
\end{quote}

\begin{quote}
\emph{kotta iṇṭlō mēmu sāmānu amarcāmu.}
\end{quote}

\begin{quote}
We set up our things in the new house.
\end{quote}

Past (I …): \te{అమర్చాను}, \emph{amarcānu}.

Putting things where they belong.

\begin{cousinweb}
\textbf{\te{అమర్చు}} (\emph{amarcu}) is native Dravidian: DEDR 162 lists Telugu \emph{amaru}, \textquotedblleft{}to suit, to be fit\textquotedblright{}, and \emph{amar(u)cu}, \textquotedblleft{}to fit, to arrange\textquotedblright{}, with Tamil \emph{amar}.
\end{cousinweb}

\begin{grammarlens}[title={What you've built}]
One more word for describing what you lived and felt.
\end{grammarlens}

\subsection*{Your turn}
\begin{itemize}
\raggedright
  \item \emph{Say it:} \emph{amarcu}
  \item \emph{Say it:} \emph{amarcu}, once more
  \item \emph{Say it:} \emph{kotta iṇṭlō mēmu sāmānu amarcāmu.} — We set up our things in the new house.
  \item \emph{Recall:} say the Telugu for to sift, then the Telugu for tempering, then say \emph{amarcu} again
\end{itemize}

\begin{tcolorbox}[breakable,colback=teal!4,colframe=teal!35!black,title={Before you move on}]
What does \te{అమర్చు} mean? (\textquotedblleft{}To set up; to arrange\textquotedblright{}.) Say it once more.
\end{tcolorbox}
\section[\texorpdfstring{\te{పరచు} (paracu)}{పరచు (paracu)}]{\te{పరచు} (paracu) — to spread out; to lay}

\label{lesson:TE-C549-paracu}



\begin{quote}
Before the new one: say the Telugu for tempering, then the Telugu for to set up.
\end{quote}

\subsection*{You'll want to know: \te{పరచు}}
\textbf{\te{పరచు}} — \emph{paracu} — \textquotedblleft{}to spread out; to lay\textquotedblright{}.

\begin{quote}
\te{అతిథుల} \te{కోసం} \te{అమ్మ} \te{నేల} \te{మీద} \te{చాప} \te{పరిచింది}.
\end{quote}

\begin{quote}
\emph{atithula kōsaṁ amma nēla mīda cāpa paricindi.}
\end{quote}

\begin{quote}
Mum spread a mat on the floor for the guests.
\end{quote}

Past (I …): \te{పరిచాను}, \emph{paricānu}.

For mats, sheets and beds.

\begin{cousinweb}
\textbf{\te{పరచు}} (\emph{paracu}) is native Dravidian: DEDR 3949 lists Telugu \emph{paravu}, \textquotedblleft{}to spread\textquotedblright{}, and \emph{parapu}, \textquotedblleft{}spreading; bedding\textquotedblright{}.
\end{cousinweb}

\begin{grammarlens}[title={What you've built}]
One more word for describing what you lived and felt.
\end{grammarlens}

\subsection*{Your turn}
\begin{itemize}
\raggedright
  \item \emph{Say it:} \emph{paracu}
  \item \emph{Say it:} \emph{paracu}, once more
  \item \emph{Say it:} \emph{atithula kōsaṁ amma nēla mīda cāpa paricindi.} — Mum spread a mat on the floor for the guests.
  \item \emph{Recall:} say the Telugu for tempering, then the Telugu for to set up, then say \emph{paracu} again
\end{itemize}

\begin{tcolorbox}[breakable,colback=teal!4,colframe=teal!35!black,title={Before you move on}]
What does \te{పరచు} mean? (\textquotedblleft{}To spread out; to lay\textquotedblright{}.) Say it once more.
\end{tcolorbox}
\section[\texorpdfstring{\te{వెలిగించు} (veligiñcu)}{వెలిగించు (veligiñcu)}]{\te{వెలిగించు} (veligiñcu) — to kindle; to light}

\label{lesson:TE-C549-veligincu}



\begin{quote}
Before the new one: say the Telugu for to set up, then the Telugu for to spread out.
\end{quote}

\subsection*{You'll want to know: \te{వెలిగించు}}
\textbf{\te{వెలిగించు}} — \emph{veligiñcu} — \textquotedblleft{}to kindle; to light\textquotedblright{}.

\begin{quote}
\te{సాయంత్రం} \te{అమ్మ} \te{దీపం} \te{వెలిగించింది}.
\end{quote}

\begin{quote}
\emph{sāyantraṁ amma dīpaṁ veligiñcindi.}
\end{quote}

\begin{quote}
In the evening mum lit the lamp.
\end{quote}

Past (I …): \te{వెలిగించాను}, \emph{veligiñcānu}.

For lamps, candles and stoves.

\begin{cousinweb}
\textbf{\te{వెలిగించు}} (\emph{veligiñcu}) is native Dravidian: DEDR 5496 lists Telugu \emph{velũgu}, \textquotedblleft{}to shine, to give light, to burn\textquotedblright{}, and its causative \emph{velĩgincu}.
\end{cousinweb}

\begin{grammarlens}[title={What you've built}]
One more word for describing what you lived and felt.
\end{grammarlens}

\subsection*{Your turn}
\begin{itemize}
\raggedright
  \item \emph{Say it:} \emph{veligiñcu}
  \item \emph{Say it:} \emph{veligiñcu}, once more
  \item \emph{Say it:} \emph{sāyantraṁ amma dīpaṁ veligiñcindi.} — In the evening mum lit the lamp.
  \item \emph{Recall:} say the Telugu for to set up, then the Telugu for to spread out, then say \emph{veligiñcu} again
\end{itemize}

\begin{tcolorbox}[breakable,colback=teal!4,colframe=teal!35!black,title={Before you move on}]
What does \te{వెలిగించు} mean? (\textquotedblleft{}To kindle; to light\textquotedblright{}.) Say it once more.
\end{tcolorbox}
\section[\texorpdfstring{\te{చిలుకు} (ciluku)}{చిలుకు (ciluku)}]{\te{చిలుకు} (ciluku) — to churn}

\label{lesson:TE-C549-ciluku}



\begin{quote}
Before the new one: say the Telugu for to spread out, then the Telugu for to kindle.
\end{quote}

\subsection*{You'll want to know: \te{చిలుకు}}
\textbf{\te{చిలుకు}} — \emph{ciluku} — \textquotedblleft{}to churn\textquotedblright{}.

\begin{quote}
\te{అమ్మమ్మ} \te{పెరుగు} \te{చిలికి} \te{వెన్న} \te{తీసింది}.
\end{quote}

\begin{quote}
\emph{ammamma perugu ciliki venna tīsindi.}
\end{quote}

\begin{quote}
Grandma churned the curd and took out the butter.
\end{quote}

Past (I …): \te{చిలికాను}, \emph{cilikānu}.

Churning curd for butter or buttermilk.

\begin{cousinweb}
\textbf{\te{చిలుకు}} (\emph{ciluku}) is native Dravidian: DEDR 2570 lists Telugu \emph{ciluku}, \textquotedblleft{}to churn\textquotedblright{}, with Tamil \emph{ciluppu}, \textquotedblleft{}to churn, to stir\textquotedblright{}.
\end{cousinweb}

\begin{grammarlens}[title={What you've built}]
One more word for describing what you lived and felt.
\end{grammarlens}

\subsection*{Your turn}
\begin{itemize}
\raggedright
  \item \emph{Say it:} \emph{ciluku}
  \item \emph{Say it:} \emph{ciluku}, once more
  \item \emph{Say it:} \emph{ammamma perugu ciliki venna tīsindi.} — Grandma churned the curd and took out the butter.
  \item \emph{Recall:} say the Telugu for to spread out, then the Telugu for to kindle, then say \emph{ciluku} again
\end{itemize}

\begin{tcolorbox}[breakable,colback=teal!4,colframe=teal!35!black,title={Before you move on}]
What does \te{చిలుకు} mean? (\textquotedblleft{}To churn\textquotedblright{}.) Say it once more.
\end{tcolorbox}
\section[\texorpdfstring{\te{అలసిపోవు} (alasipōvu)}{అలసిపోవు (alasipōvu)}]{\te{అలసిపోవు} (alasipōvu) — to get tired; to wear out}

\label{lesson:TE-C549-alasipovu}



\begin{quote}
Before the new one: say the Telugu for to kindle, then the Telugu for to churn.
\end{quote}

\subsection*{You'll want to know: \te{అలసిపోవు}}
\textbf{\te{అలసిపోవు}} — \emph{alasipōvu} — \textquotedblleft{}to get tired; to wear out\textquotedblright{}.

\begin{quote}
\te{ఇల్లు} \te{సర్దిన} \te{తరువాత} \te{మేము} \te{బాగా} \te{అలసిపోయాము}.
\end{quote}

\begin{quote}
\emph{illu sardina taruvāta mēmu bāgā alasipōyāmu.}
\end{quote}

\begin{quote}
After tidying the house we were very tired.
\end{quote}

Past (I …): \te{అలసిపోయాను}, \emph{alasipōyānu}.

Its noun \te{అలసట} is tiredness.

\begin{cousinweb}
\textbf{\te{అలసిపోవు}} (\emph{alasipōvu}) is native Dravidian: DEDR 236 lists Telugu \emph{alayu}, \textquotedblleft{}to be tired\textquotedblright{}, and \emph{alasaṭa}, \textquotedblleft{}fatigue\textquotedblright{}; the verb ends in \emph{pōvu}, \textquotedblleft{}to go\textquotedblright{} (DEDR 4572).
\end{cousinweb}

\begin{grammarlens}[title={What you've built}]
One more word for describing what you lived and felt.
\end{grammarlens}

\subsection*{Your turn}
\begin{itemize}
\raggedright
  \item \emph{Say it:} \emph{alasipōvu}
  \item \emph{Say it:} \emph{alasipōvu}, once more
  \item \emph{Say it:} \emph{illu sardina taruvāta mēmu bāgā alasipōyāmu.} — After tidying the house we were very tired.
  \item \emph{Recall:} say the Telugu for to kindle, then the Telugu for to churn, then say \emph{alasipōvu} again
\end{itemize}

\begin{tcolorbox}[breakable,colback=teal!4,colframe=teal!35!black,title={Before you move on}]
What does \te{అలసిపోవు} mean? (\textquotedblleft{}To get tired; to wear out\textquotedblright{}.) Say it once more.
\end{tcolorbox}
\section[(dialogue)]{Making a home, again}

\label{lesson:TE-R549-review-542-549}



\begin{quote}
Say the Telugu for the art of painting and to portray.
\end{quote}

\subsection*{Your turn}
Say these aloud:

\begin{itemize}
\raggedright
  \item the art of painting, to portray, a sculpture, a sculptor, a masterpiece
  \item a concert, applause, to congratulate, to savour, a stick dance
  \item a sportsman, an opponent, to clash, to prevail, running
  \item to study closely, a workout, to aim, to shove, a bet
  \item the eaves, a low wooden stool, a dug well, a small water pot, a cockroach
  \item to cleanse, to wash clothes, to hang to dry, to shake out, to hang up
  \item to boil milk, to scorch, to chop, to sift, tempering
  \item to set up, to spread out, to kindle, to churn, to get tired
\end{itemize}

\begin{tcolorbox}[breakable,colback=teal!4,colframe=teal!35!black,title={Before you move on}]
The art of painting? (\textbf{\te{చిత్రలేఖనం}}, \emph{citralēkhanaṁ}.) To portray? (\textbf{\te{చిత్రించు}}, \emph{citriñcu}.) A sculpture? (\textbf{\te{శిల్పం}}, \emph{śilpaṁ}.) A sculptor? (\textbf{\te{శిల్పి}}, \emph{śilpi}.) A masterpiece? (\textbf{\te{కళాఖండం}}, \emph{kaḷākhaṇḍaṁ}.) A concert? (\textbf{\te{కచేరీ}}, \emph{kacērī}.) Applause? (\textbf{\te{చప్పట్లు}}, \emph{cappaṭlu}.) To congratulate? (\textbf{\te{అభినందించు}}, \emph{abhinandiñcu}.) To savour? (\textbf{\te{ఆస్వాదించు}}, \emph{āsvādiñcu}.) A stick dance? (\textbf{\te{కోలాటం}}, \emph{kōlāṭaṁ}.) A sportsman? (\textbf{\te{క్రీడాకారుడు}}, \emph{krīḍākāruḍu}.) An opponent? (\textbf{\te{ప్రత్యర్థి}}, \emph{pratyarthi}.) To clash? (\textbf{\te{తలపడు}}, \emph{talapaḍu}.) To prevail? (\textbf{\te{నెగ్గు}}, \emph{neggu}.) Running? (\textbf{\te{పరుగు}}, \emph{parugu}.) To study closely? (\textbf{\te{అభ్యసించు}}, \emph{abhyasiñcu}.) A workout? (\textbf{\te{కసరత్తు}}, \emph{kasarattu}.) To aim? (\textbf{\te{గురిపెట్టు}}, \emph{guripeṭṭu}.) To shove? (\textbf{\te{తోయు}}, \emph{tōyu}.) A bet? (\textbf{\te{పందెం}}, \emph{pandeṁ}.) The eaves? (\textbf{\te{చూరు}}, \emph{cūru}.) A low wooden stool? (\textbf{\te{పీట}}, \emph{pīṭa}.) A dug well? (\textbf{\te{బావి}}, \emph{bāvi}.) A small water pot? (\textbf{\te{చెంబు}}, \emph{cembu}.) A cockroach? (\textbf{\te{బొద్దింక}}, \emph{boddiṅka}.) To cleanse? (\textbf{\te{శుభ్రపరచు}}, \emph{śubhraparacu}.) To wash clothes? (\textbf{\te{ఉతుకు}}, \emph{utuku}.) To hang to dry? (\textbf{\te{ఆరవేయు}}, \emph{āravēyu}.) To shake out? (\textbf{\te{దులుపు}}, \emph{dulupu}.) To hang up? (\textbf{\te{తగిలించు}}, \emph{tagiliñcu}.) To boil milk? (\textbf{\te{కాచు}}, \emph{kācu}.) To scorch? (\textbf{\te{మాడ్చు}}, \emph{māḍcu}.) To chop? (\textbf{\te{తరుగు}}, \emph{tarugu}.) To sift? (\textbf{\te{జల్లించు}}, \emph{jalliñcu}.) Tempering? (\textbf{\te{తాలింపు}}, \emph{tālimpu}.) To set up? (\textbf{\te{అమర్చు}}, \emph{amarcu}.) To spread out? (\textbf{\te{పరచు}}, \emph{paracu}.) To kindle? (\textbf{\te{వెలిగించు}}, \emph{veligiñcu}.) To churn? (\textbf{\te{చిలుకు}}, \emph{ciluku}.) To get tired? (\textbf{\te{అలసిపోవు}}, \emph{alasipōvu}.)
\end{tcolorbox}
\section[(dialogue)]{Making a Home — once more}

\label{lesson:TE-R549-review-final-549}



\begin{quote}
Say the Telugu for to set up and to spread out.
\end{quote}

\subsection*{Your turn}
\begin{itemize}
\raggedright
  \item \emph{Say it:} to set up, to spread out, to kindle, to churn, to get tired
\end{itemize}

\begin{tcolorbox}[breakable,colback=teal!4,colframe=teal!35!black,title={Before you move on}]
To set up? (\textbf{\te{అమర్చు}}, \emph{amarcu}.) To spread out? (\textbf{\te{పరచు}}, \emph{paracu}.) To kindle? (\textbf{\te{వెలిగించు}}, \emph{veligiñcu}.) To churn? (\textbf{\te{చిలుకు}}, \emph{ciluku}.) To get tired? (\textbf{\te{అలసిపోవు}}, \emph{alasipōvu}.)
\end{tcolorbox}
